\chapter{Matrix Algebra and Inverses}

A matrix began in Chapter~2 as bookkeeping for a linear system. This chapter treats matrices as objects of their own: they can be added, scaled and, most importantly, multiplied. Matrix multiplication is the operation that makes the rest of the subject work. We then ask when a matrix can be ``divided by'', that is, when it has an inverse, and we connect invertibility to the elimination of Chapter~2 through elementary matrices.

\section{Matrices and their operations}

Write $\R^{m\times n}$ for the set of real $m\times n$ matrices. The entry in row $i$ and column $j$ of $A$ is $a_{ij}$ (also written $(A)_{ij}$). A matrix with $m=n$ is \emph{square}. The matrix with all entries $0$ is $0$; the \emph{identity matrix} $I_n$ is the $n\times n$ matrix with ones on the diagonal and zeros elsewhere.

\begin{definition}[Addition and scalar multiplication]
For $A,B\in\R^{m\times n}$ and $c\in\R$, $(A+B)_{ij}=a_{ij}+b_{ij}$ and $(cA)_{ij}=c\,a_{ij}$.
\end{definition}

\begin{definition}[Matrix product]
Let $A$ be $m\times n$ and $B$ be $n\times p$. The product $AB$ is the $m\times p$ matrix with
\[ (AB)_{ij}=\sum_{k=1}^n a_{ik}b_{kj}. \]
The product is defined only when the number of columns of $A$ equals the number of rows of $B$.
\end{definition}

The formula says that $(AB)_{ij}$ is the dot product of row $i$ of $A$ with column $j$ of $B$. There are three other equivalent descriptions, and fluency with all four is a mark of understanding.

\begin{proposition}[Four views of $AB$]\label{prop:fourviews}
Let $A$ have rows $\mathbf{r}_1^\top,\dots,\mathbf{r}_m^\top$ and columns $\mathbf{a}_1,\dots,\mathbf{a}_n$, and let $B$ have columns $\mathbf{b}_1,\dots,\mathbf{b}_p$ and rows $\mathbf{s}_1^\top,\dots,\mathbf{s}_n^\top$.
\begin{enumerate}[label=(\roman*)]
  \item \textbf{Entry view.} $(AB)_{ij}=\mathbf{r}_i\cdot\mathbf{b}_j$.
  \item \textbf{Column view.} The $j$-th column of $AB$ is $A\mathbf{b}_j$, a combination of the columns of $A$ with weights from $\mathbf{b}_j$.
  \item \textbf{Row view.} The $i$-th row of $AB$ is $\mathbf{r}_i^\top B$, a combination of the rows of $B$.
  \item \textbf{Outer-product view.} $AB=\sum_{k=1}^n\mathbf{a}_k\mathbf{s}_k^\top$, a sum of $n$ rank-one matrices.
\end{enumerate}
\end{proposition}

\begin{proof}
(i) is the definition. (ii) and (iii) restate (i) for a fixed column or row, using $A\mathbf{x}=x_1\mathbf{a}_1+\dots+x_n\mathbf{a}_n$. For (iv), the $(i,j)$ entry of $\mathbf{a}_k\mathbf{s}_k^\top$ is $a_{ik}b_{kj}$; summing over $k$ gives the definition.
\end{proof}

\begin{intuition}
The column view says that multiplying by $A$ builds new vectors out of the columns of $A$. This is the view that matters for linear maps in Chapter~9: the columns of $A$ are the images of the standard basis vectors. The outer-product view is the basis of low-rank approximation: a matrix of rank $r$ is a sum of $r$ outer products, and the SVD of Chapter~17 finds the best such sums.
\end{intuition}

\begin{example}
\[
\begin{pmatrix}1&2\\3&4\end{pmatrix}\begin{pmatrix}0&1\\1&0\end{pmatrix}=\begin{pmatrix}2&1\\4&3\end{pmatrix},\qquad
\begin{pmatrix}0&1\\1&0\end{pmatrix}\begin{pmatrix}1&2\\3&4\end{pmatrix}=\begin{pmatrix}3&4\\1&2\end{pmatrix}.
\]
Multiplying on the right by the swap matrix swaps the \emph{columns}; on the left it swaps the \emph{rows}. The two products differ.
\end{example}

\begin{theorem}[Rules of matrix algebra]\label{thm:matalg}
Whenever the sizes make the expressions defined:
\begin{enumerate}[label=(\roman*)]
  \item $(AB)C=A(BC)$ (associativity);
  \item $A(B+C)=AB+AC$ and $(A+B)C=AC+BC$ (distributivity);
  \item $c(AB)=(cA)B=A(cB)$;
  \item $I_mA=A=AI_n$ for $A\in\R^{m\times n}$.
\end{enumerate}
\end{theorem}

\begin{proof}
For (i), the $(i,l)$ entry of $(AB)C$ is $\sum_j\big(\sum_k a_{ik}b_{kj}\big)c_{jl}=\sum_{j,k}a_{ik}b_{kj}c_{jl}$, and the same double sum arises for $A(BC)$ after exchanging the order of summation. The others are equally direct.
\end{proof}

\begin{pitfall}{matrix multiplication is not commutative}
In general $AB\ne BA$; one product may exist while the other does not (for $A$ of size $2\times3$ and $B$ of size $3\times 4$). Consequences that fail for matrices but hold for numbers: $(A+B)^2\ne A^2+2AB+B^2$ in general, $AB=0$ does not imply $A=0$ or $B=0$, and $AB=AC$ does not imply $B=C$. For example $\begin{pmatrix}1&0\\0&0\end{pmatrix}\begin{pmatrix}0&0\\0&1\end{pmatrix}=0$.
\end{pitfall}

\subsection*{Powers and the transpose}

For a square $A$ define $A^0=I$, $A^k=A\cdots A$ ($k$ factors). Then $A^{j+k}=A^jA^k$, but $(AB)^k\ne A^kB^k$ unless $A,B$ commute.

\begin{definition}[Transpose]
The \emph{transpose} $A\T$ of $A\in\R^{m\times n}$ is the $n\times m$ matrix with $(A\T)_{ij}=a_{ji}$. A square matrix is \emph{symmetric} if $A\T=A$ and \emph{skew-symmetric} if $A\T=-A$.
\end{definition}

\begin{proposition}[Transpose rules]\label{prop:transpose}
$(A\T)\T=A$, $(A+B)\T=A\T+B\T$, $(cA)\T=cA\T$, and $(AB)\T=B\T A\T$ (the order reverses).
\end{proposition}

\begin{proof}
For the product: $\big((AB)\T\big)_{ij}=(AB)_{ji}=\sum_k a_{jk}b_{ki}=\sum_k(B\T)_{ik}(A\T)_{kj}=(B\T A\T)_{ij}$.
\end{proof}

\begin{proposition}\label{prop:gram}
For every real matrix $A$ the matrices $A\T A$ and $AA\T$ are symmetric and positive semidefinite in the sense that $\mathbf{x}\T A\T A\mathbf{x}=\norm{A\mathbf{x}}^2\ge0$. Every symmetric matrix is the sum of a symmetric and a skew part in exactly one way: $A=\tfrac12(A+A\T)+\tfrac12(A-A\T)$.
\end{proposition}

\begin{proof}
$(A\T A)\T=A\T(A\T)\T=A\T A$. Also $\mathbf{x}\T A\T A\mathbf{x}=(A\mathbf{x})\T(A\mathbf{x})=\norm{A\mathbf{x}}^2$. The decomposition is checked directly, and is unique because a matrix that is both symmetric and skew is zero.
\end{proof}

The matrices $A\T A$ (the \emph{Gram matrix}) appear throughout statistics, least squares and the SVD.

\section{Block matrices}

A matrix can be cut into rectangular blocks by horizontal and vertical lines, and blocks multiply like entries as long as the dimensions fit:
\[
\begin{pmatrix}A&B\\C&D\end{pmatrix}\begin{pmatrix}E&F\\G&H\end{pmatrix}=\begin{pmatrix}AE+BG&AF+BH\\CE+DG&CF+DH\end{pmatrix}.
\]
This is Proposition~\ref{prop:fourviews} applied with blocks in place of entries. It is the tool behind fast algorithms (Strassen's method multiplies two $n\times n$ matrices in $O(n^{2.81})$ operations by a clever block identity) and behind the Schur complement of Chapter~16.

\section{Inverse matrices}

\begin{definition}[Invertible matrix]
A square matrix $A\in\R^{n\times n}$ is \emph{invertible} (or \emph{nonsingular}) if there is a matrix $B$ with $AB=BA=I_n$. Otherwise $A$ is \emph{singular}.
\end{definition}

\begin{proposition}[Uniqueness and rules]\label{prop:invrules}
If an inverse exists it is unique, and is written $A^{-1}$. If $A,B$ are invertible then $AB$ is invertible with $(AB)^{-1}=B^{-1}A^{-1}$, and $(A\T)^{-1}=(A^{-1})\T$.
\end{proposition}

\begin{proof}
If $B$ and $C$ are both inverses then $B=BI=B(AC)=(BA)C=IC=C$. Next, $(AB)(B^{-1}A^{-1})=A(BB^{-1})A^{-1}=I$ and likewise on the other side. Finally take transposes of $AA^{-1}=I$ and use Proposition~\ref{prop:transpose}.
\end{proof}

\begin{example}[$2\times2$ inverse]
For $A=\begin{pmatrix}a&b\\c&d\end{pmatrix}$ with $ad-bc\ne0$,
\[ A^{-1}=\frac{1}{ad-bc}\begin{pmatrix}d&-b\\-c&a\end{pmatrix}, \]
as one checks by multiplying. The quantity $ad-bc$ is the determinant (Chapter~5); if it vanishes the matrix is singular.
\end{example}

\begin{theorem}[The invertible matrix theorem, first version]\label{thm:imt1}
For a square matrix $A\in\R^{n\times n}$ the following are equivalent.
\begin{enumerate}[label=(\alph*)]
  \item $A$ is invertible.
  \item $A\mathbf{x}=\mathbf{b}$ has a solution for every $\mathbf{b}\in\R^n$.
  \item $A\mathbf{x}=\mathbf{0}$ has only the trivial solution $\mathbf{x}=\mathbf{0}$.
  \item The reduced row echelon form of $A$ is $I_n$.
  \item $A$ has $n$ pivots.
\end{enumerate}
In that case $A\mathbf{x}=\mathbf{b}$ has the unique solution $\mathbf{x}=A^{-1}\mathbf{b}$.
\end{theorem}

\begin{proof}
(a)$\Rightarrow$(b),(c): take $\mathbf{x}=A^{-1}\mathbf{b}$; and if $A\mathbf{x}=\mathbf{0}$ then $\mathbf{x}=A^{-1}A\mathbf{x}=\mathbf{0}$. (c)$\Rightarrow$(e): no free variables means every column is a pivot column, so there are $n$ pivots. (e)$\Leftrightarrow$(d): a square RREF with $n$ pivots has a pivot in every row and column, hence is $I_n$. (d)$\Rightarrow$(b): reduce $[A\mid\mathbf{b}]$; the left block becomes $I_n$ so the system is consistent. (b)$\Rightarrow$(a): solve $A\mathbf{x}_j=\mathbf{e}_j$ for $j=1,\dots,n$ and let $B=[\mathbf{x}_1\cdots\mathbf{x}_n]$, so $AB=I$. If $B\mathbf{y}=\mathbf{0}$ then $\mathbf{y}=AB\mathbf{y}=\mathbf{0}$, so $B$ satisfies (c). The implications (c)$\Rightarrow$(e)$\Rightarrow$(d)$\Rightarrow$(b), which were proved for an arbitrary square matrix, then show that $B$ satisfies (b); in particular some $C$ has $BC=I$. Now $A=A(BC)=(AB)C=C$, so $BA=BC=I$ and $B=A^{-1}$.
\end{proof}

A useful consequence: \emph{for square matrices a one-sided inverse is automatically two-sided}. If $AB=I$ then $BA=I$.

\section{Elementary matrices and the inverse algorithm}

Each row operation is itself a matrix multiplication.

\begin{definition}[Elementary matrix]
An \emph{elementary matrix} is obtained from $I_n$ by one elementary row operation. Left multiplication by it performs that operation: if $E$ is obtained from $I$ by an operation, then $EA$ is obtained from $A$ by the same operation.
\end{definition}

\begin{example}
For $3\times 3$: $E_{21}(-2)=\begin{pmatrix}1&0&0\\-2&1&0\\0&0&1\end{pmatrix}$ subtracts twice row $1$ from row $2$; its inverse $E_{21}(2)$ adds it back. The swap $P_{12}$ and the scaling $D=\diag(1,c,1)$ are the other two kinds.
\end{example}

\begin{theorem}[Elementary matrices are invertible]
Every elementary matrix is invertible, and its inverse is an elementary matrix of the same type.
\end{theorem}

\begin{proof}
Each operation is undone by another operation (swap again; scale by $1/c$; add the multiple with the opposite sign), and the corresponding elementary matrix performs it.
\end{proof}

\begin{theorem}[Invertible = product of elementary matrices]\label{thm:elemprod}
A square matrix is invertible if and only if it is a product of elementary matrices.
\end{theorem}

\begin{proof}
If $A$ is invertible, Theorem~\ref{thm:imt1} shows that a sequence of row operations reduces $A$ to $I$: $E_k\cdots E_1A=I$. Hence $A=E_1^{-1}\cdots E_k^{-1}$, a product of elementary matrices. Conversely, products of invertible matrices are invertible.
\end{proof}

\subsection*{Computing the inverse}

If $E_k\cdots E_1A=I$ then $A^{-1}=E_k\cdots E_1$, which is the result of applying the same row operations to $I$. Therefore:

\begin{theorem}[Gauss--Jordan inverse algorithm]
Row reduce the $n\times 2n$ matrix $[A\mid I]$. If the left block becomes $I$, the right block is $A^{-1}$; if the left block cannot reach $I$, then $A$ is singular.
\end{theorem}

\begin{example}
For $A=\begin{pmatrix}1&2\\3&7\end{pmatrix}$:
\[
\begin{pmatrix}1&2&1&0\\3&7&0&1\end{pmatrix}\xrightarrow{R_2-3R_1}\begin{pmatrix}1&2&1&0\\0&1&-3&1\end{pmatrix}\xrightarrow{R_1-2R_2}\begin{pmatrix}1&0&7&-2\\0&1&-3&1\end{pmatrix},
\]
so $A^{-1}=\begin{pmatrix}7&-2\\-3&1\end{pmatrix}$.
\end{example}

\begin{pitfall}{computing inverses to solve systems}
The formula $\mathbf{x}=A^{-1}\mathbf{b}$ is for reasoning, not computing. Forming $A^{-1}$ costs about three times as much as one elimination, is less accurate, and destroys sparsity (the inverse of a sparse matrix is usually dense). In code, solve the system directly (\texttt{solve(A, b)}) and never write \texttt{inv(A) @ b}.
\end{pitfall}

\begin{interview}{inverse of a sum and low-rank updates}
The identity $(I+uv\T)^{-1}=I-\dfrac{uv\T}{1+v\T u}$ (valid when $1+v\T u\ne0$) is the \emph{Sherman--Morrison formula}. It updates an inverse after a rank-one change in $O(n^2)$ instead of $O(n^3)$ operations, and underlies recursive least squares and the Kalman filter. Be ready to verify it by multiplying out: $(I+uv\T)\big(I-\frac{uv\T}{1+v\T u}\big)=I+uv\T-\frac{uv\T+u(v\T u)v\T}{1+v\T u}=I$.
\end{interview}

\section{Special families of matrices}

\begin{center}
\begin{tabular}{lll}
\toprule
Family & Definition & Closed under product? \\
\midrule
Diagonal & $a_{ij}=0$ for $i\ne j$ & yes \\
Upper triangular & $a_{ij}=0$ for $i>j$ & yes \\
Lower triangular & $a_{ij}=0$ for $i<j$ & yes \\
Symmetric & $A\T=A$ & no \\
Permutation & one $1$ in each row and column & yes \\
Orthogonal & $A\T A=I$ & yes (Chapter~11) \\
\bottomrule
\end{tabular}
\end{center}

An upper triangular matrix is invertible if and only if all diagonal entries are nonzero, and its inverse is again upper triangular. A \emph{permutation matrix} reorders coordinates and satisfies $P^{-1}=P\T$.

\begin{summarybox}
\begin{itemize}
  \item Matrix multiplication has four equivalent descriptions (entry, column, row, outer product); it is associative and distributive, but not commutative.
  \item $(AB)\T=B\T A\T$ and $(AB)^{-1}=B^{-1}A^{-1}$: order reverses.
  \item A square matrix is invertible exactly when it has $n$ pivots, when $A\mathbf{x}=\mathbf{0}$ has only the zero solution, and when it is a product of elementary matrices.
  \item $[A\mid I]\to[I\mid A^{-1}]$ computes inverses, but solving systems directly is cheaper and safer.
  \item The Gram matrices $A\T A$ and $AA\T$ are symmetric and positive semidefinite.
\end{itemize}
\end{summarybox}

\section{Exercises}

\begin{exercise}
Compute $AB$ and $BA$ for $A=\begin{pmatrix}1&2&0\\0&1&3\end{pmatrix}$ and $B=\begin{pmatrix}1&0\\2&1\\0&4\end{pmatrix}$, and explain why the sizes of the two products differ.
\end{exercise}

\begin{exercise}
Find $2\times 2$ matrices $A,B$ with $AB=0$ but $BA\ne0$.
\end{exercise}

\begin{exercise}
Prove $(ABC)\T=C\T B\T A\T$ and $(ABC)^{-1}=C^{-1}B^{-1}A^{-1}$.
\end{exercise}

\begin{exercise}
Show that if $A^2=A$ (an \emph{idempotent} matrix) and $A$ is invertible, then $A=I$.
\end{exercise}

\begin{exercise}
Invert $\begin{pmatrix}1&0&2\\2&-1&3\\4&1&8\end{pmatrix}$ by the Gauss--Jordan algorithm and check by multiplication.
\end{exercise}

\begin{exercise}
Let $N$ be square with $N^k=0$ for some $k$ (a \emph{nilpotent} matrix). Prove $I-N$ is invertible with $(I-N)^{-1}=I+N+\dots+N^{k-1}$.
\end{exercise}

\begin{exercise}
Prove that every symmetric $2\times2$ matrix $\begin{pmatrix}a&b\\b&c\end{pmatrix}$ with $ac-b^2>0$ and $a>0$ satisfies $\mathbf{x}\T A\mathbf{x}>0$ for all $\mathbf{x}\ne\mathbf{0}$. (Complete the square.)
\end{exercise}

\begin{exercise}
Verify the Sherman--Morrison formula for $u=(1,0)\T$, $v=(1,2)\T$ by computing both sides explicitly.
\end{exercise}
